\scanpage{66}
% archive sheet number: 67
\deleted{transforme hermitiens en hermitiens)}. Il s'ensuit bien que $\overset{n}{\Lambda}h$ \deleted{est}, carré d'un hermitien, est hermitien \uline{positif}.

Enfin, la formule (11) est vraie si $\alpha$ est entier, elle est aussi vraie si $\alpha=\frac1n$ ($n$ entier $>0$) car on a $\bigl(\overset{n}{\Lambda}h^{1/n}\bigr)^n=\overset{n}{\Lambda}\bigl((h^{1/n})^n\bigr)=\overset{n}{\Lambda}h$, \deleted{et on voit que} \inserted{$\overset{n}{\Lambda}h^{1/n}$ est un} \uline{hermit.} \uline{positif}, il en résulte que c'est bien la racine $n^{\text{ième}}$ de $\overset{n}{\Lambda}h$. \deleted{Par suite, que (12) est \deleted{\unsure{vrai}} vraie pour le \unsure{fractionnaire},} En combinant on voit que (12) est vraie pour $\alpha$ fractionnaire.

D'autre part, si $H$ est un \deleted{\unsure{hermitien}}\inserted{\uline{opérateur}} positif dans un espace de Hilbert, on sait que $\alpha\longrightarrow H^\alpha$ est une fonction continue de $\alpha>0$. \deleted{\unsure{à}} Donc, l'application $U\longrightarrow\overset{n}{\Lambda}u$ de $L(E)$ dans $L(\overset{n}{\Lambda}^{(2)}E)$ étant aussi continue, on voit que les fonctions $(\overset{n}{\Lambda}h)^\alpha$ et $\overset{n}{\Lambda}(h^{1/\alpha})$ sont fonctions continues de $\alpha$, \deleted{et} comme elles sont égales pour des valeurs rationnelles, elles sont égales partout, ce qui achève d'établir (12). Résumons l'essentiel de ces résultats

\uline{Proposition} 1 Soient $E$ et $F$ deux espaces de Hilbert, soient $u_i$ \deleted{\unsure{$(u_1,\ldots,u_n)$}} \inserted{$(i=1,\ldots n)$} des opérateurs linéaires continus de $E$ dans $F$, alors $u_1\wedge\ldots\wedge u_n$ est continu pour les structures préhilbertiennes de $\overset{n}{\Lambda}^{a}E$ et $\overset{n}{\Lambda}^{a}F$, et se prolonge donc par continuité en une application linéaire continue de $\overset{n}{\Lambda}^{(1)}E$ dans $\overset{n}{\Lambda}^{(2)}F$, \deleted{et} c. à d. un $u_1\wedge\ldots\wedge u_n$, \deleted{et satisfait à} de norme $\leq\lVert u_1\rVert\ldots\lVert u_n\rVert$.

L'application $u\longrightarrow\overset{n}{\Lambda}u$ de $L(E)$ dans $L(\overset{n}{\Lambda}^{(2)}E)$ est multiplicative, commute à l'opération $*$, transforme op. hermitiens en op. hermitiens et opérateurs positifs en
